\subsection{等参有限元}
\label{sec:ch10-isoparametric}
\setcounter{thmcount}{0}
\setcounter{equation}{0}

使用等参有限元(第~\MBUnresolvedReference{sec:ch04-isoparametric-elements}{4.7} 节)同样会造成变分罪行: 可以把它理解为违反式~\eqref{eq:ch10-conforming-inclusion}, 也可以理解为使用了修正的双线性形式. 不过, 这种方法比上一节的方法更容易估计. 用多项式插值边界条件时, $v\in V_h$ 在 $\partial\Omega$ 上可能显著偏离零. Lobatto 插值点保证误差以能相互抵消不利影响的方式振荡, 从而减小这种影响. 另一方面, 等参有限元通过从某个区域 $\Omega_h$ 出发的分片多项式映射变换而来, 因而不再是 $\Omega$ 的坐标多项式. 所以它们能在逐点意义下更精确地匹配 Dirichlet 边界条件.

先给出模型问题~\eqref{eq:ch10-dirichlet-model} 的表述. 回顾我们有 $\Omega$ 的多面体逼近 $\Omega_h$ 和映射 $F^h$, 使得 $F^h(\Omega_h)$ 紧密逼近 $\Omega$. 假设等参映射满足第~\MBUnresolvedReference{sec:ch04-isoparametric-elements}{4.7} 节的性质 1--3. 特别地, 假设存在辅助映射 $F:\Omega_h\to\Omega$, 且映射的每个分量都满足 $F_i^h=I^hF_i$. 还假设 $F$ 也满足第~\MBUnresolvedReference{sec:ch04-isoparametric-elements}{4.7} 节的条件 1 和 3. 构造这样的 $F$ 并非易事, 但 Lenoir (1986) 给出了构造. 定义
\[
\MathBookFormulaID{formula-ch10-isoparametric-bilinear-form}
a_h(v,w)=\int_{F^h(\Omega_h)}\nabla v(x)\cdot\nabla w(x)\,dx.
\]

\MathBookSourcePage{297}
为便于分析, 需要映射 $\Phi^h:\Omega\to F^h(\Omega_h)$, 定义为 $\Phi^h(x)=F^h(F^{-1}(x))$. 可写成
\[
\MathBookFormulaID{formula-ch10-isoparametric-pulled-form}
a_h(v,w)=\int_\Omega J_{\Phi^h}(x)^{-t}\nabla\widehat v(x)\cdot
J_{\Phi^h}(x)^{-t}\nabla\widehat w(x)\det J_{\Phi^h}(x)\,dx,
\]
其中, 对定义在 $F^h(\Omega_h)$ 上的任意函数 $v$, 令
\MBSourceItem{1}
\begin{equation}
\MathBookFormulaID{formula-ch10-isoparametric-pullback}
\label{eq:ch10-isoparametric-pullback}
\widehat v(x):=v(\Phi^h(x)),
\end{equation}
且 $J_{\Phi^h}$ 表示 $\Phi^h$ 的 Jacobian. 关键在于
\MBSourceItem{2}
\begin{equation}
\MathBookFormulaID{formula-ch10-isoparametric-jacobian-error}
\label{eq:ch10-isoparametric-jacobian-error}
J_{\Phi^h}-I=O(h^{k-1}).
\end{equation}

令 $V_h$ 如式~\MBUnresolvedReference{eq:ch04-source-4-7-2}{4.7.2} 所定义, 其中 $\widetilde\Omega=\Omega_h$, $\widetilde F=F^h$, 并在 $V_h$ 中定义 $u_h$, 使得
\MBSourceItem{3}
\begin{equation}
\MathBookFormulaID{formula-ch10-isoparametric-discrete-problem}
\label{eq:ch10-isoparametric-discrete-problem}
a_h(u_h,v)=\int_{F^h(\Omega_h)}f(x)v(x)\,dx.
\end{equation}
先推导与式~\eqref{eq:ch10-second-strang-estimate} 类似的估计. 定义 $\widehat V_h=\{\widehat v:v\in V_h\}$, 并令 $\widehat v_h$ 表示 $u$ 到 $\widehat V_h$ 上的投影:
\MBSourceItem{4}
\begin{equation}
\MathBookFormulaID{formula-ch10-isoparametric-projection}
\label{eq:ch10-isoparametric-projection}
a(\widehat v_h,\widehat w)=a(u,\widehat w)
\qquad\forall\widehat w\in\widehat V_h.
\end{equation}
注意 $\widehat V_h\subset V$, 即 $\widehat w\in\widehat V_h$ 严格满足 Dirichlet 边界条件. 对任意 $w\in V_h$,
\MBSourceItem{5}
\begin{equation}
\MathBookFormulaID{formula-ch10-isoparametric-consistency-identity}
\label{eq:ch10-isoparametric-consistency-identity}
\begin{split}
a(\widehat v_h-\widehat u_h,\widehat w)
&=a(u-\widehat u_h,\widehat w)\\
&=(f,\widehat w)-a(\widehat u_h,\widehat w)\\
&=(f,\widehat w)-\int_{F^h(\Omega_h)}fw\,dx
+a_h(u_h,w)-a(\widehat u_h,\widehat w).
\end{split}
\end{equation}
所以
\MBSourceItem{6}
\begin{equation}
\MathBookFormulaID{formula-ch10-isoparametric-abstract-estimate}
\label{eq:ch10-isoparametric-abstract-estimate}
\begin{split}
\|u-\widehat u_h\|_a
&\leq\inf_{v\in V_h}\|u-\widehat v\|_a
+\sup_{w\in V_h\setminus\{0\}}
\frac{|a_h(u_h,w)-a(\widehat u_h,\widehat w)|}{\|\widehat w\|_a}\\
&\quad+\sup_{w\in V_h\setminus\{0\}}
\frac{\left|(f,\widehat w)-\int_{F^h(\Omega_h)}fw\,dx\right|}{\|\widehat w\|_a}.
\end{split}
\end{equation}

\MathBookSourcePage{298}
上式第二项中的分子可展开为
\begin{align*}
\MathBookFormulaID{formula-ch10-isoparametric-form-expansion}
a_h(u_h,w)-a(\widehat u_h,\widehat w)
&=\int_\Omega(J_{\Phi^h}^{-t}-I)\nabla\widehat u_h\cdot
J_{\Phi^h}^{-t}\nabla\widehat w\det J_{\Phi^h}\,dx\\
&\quad+\int_\Omega\nabla\widehat u_h\cdot
\bigl((\det J_{\Phi^h})J_{\Phi^h}^{-t}-I\bigr)\nabla\widehat w\,dx.
\end{align*}
由式~\eqref{eq:ch10-isoparametric-jacobian-error},
\MBSourceItem{7}
\begin{equation}
\MathBookFormulaID{formula-ch10-isoparametric-form-error}
\label{eq:ch10-isoparametric-form-error}
\sup_{w\in V_h\setminus\{0\}}
\frac{|a_h(u_h,w)-a(\widehat u_h,\widehat w)|}{\|\widehat w\|_a}
\leq Ch^{k-1}\|\widehat u_h\|_a.
\end{equation}
含 $f$ 的项可类似估计:
\begin{align*}
\MathBookFormulaID{formula-ch10-isoparametric-load-error}
\left|(f,\widehat w)-\int_{F^h(\Omega_h)}fw\,dx\right|
&=\left|\int_\Omega\bigl(f(x)-f(\Phi^h(x))\det J_{\Phi^h}(x)\bigr)\widehat w(x)\,dx\right|\\
&\leq C\bigl(h^k\|f\|_{W_\infty^1(\Omega)}+h^{k-1}\|f\|_{L^2(\Omega)}\bigr)
\|\widehat w\|_{L^2(\Omega)}.
\end{align*}
结合式~\eqref{eq:ch10-isoparametric-abstract-estimate} 和~\eqref{eq:ch10-isoparametric-form-error}, 得到下述定理.

\MBSourceItem{8}
\begin{Theorem}
\label{thm:ch10-isoparametric-error}
设 $u$ 由问题~\eqref{eq:ch10-dirichlet-model} 定义, $u_h$ 由式~\eqref{eq:ch10-isoparametric-discrete-problem} 定义, 其中 $V_h$ 由式~\MBUnresolvedReference{eq:ch04-source-4-7-2}{4.7.2} 定义($\widetilde\Omega=\Omega_h$, $\widetilde F=F^h$). 当 $h$ 充分小时,
\[
\MathBookFormulaID{formula-ch10-isoparametric-error}
\|u-\widehat u_h\|_{H^1(\Omega)}
\leq Ch^{k-1}\bigl(\|u\|_{H^k(\Omega)}+\|f\|_{W_\infty^1(\Omega)}\bigr),
\]
其中 $\widehat u_h$ 的定义见式~\eqref{eq:ch10-isoparametric-pullback}.
\end{Theorem}
